% ── Section 1: Introduction ──────────────────────────────────────────
\section{Introduction}

The rapid adoption of Large Language Models (LLMs) in user-facing applications has fundamentally changed how software systems process natural-language input. Products such as ChatGPT, Bing Chat, GitHub Copilot, and a growing family of autonomous ``AI agents'' now handle tasks ranging from customer support and code generation to email drafting and web browsing~\cite{achiam2023gpt4}. These systems typically operate by concatenating a developer-authored \emph{system prompt}---which defines the application's behaviour---with untrusted user input, then feeding the combined text to a single language model. Because the model has no built-in mechanism to distinguish instructions from data, this architecture creates a new class of \emph{injection} vulnerability that is conceptually similar to SQL injection but far harder to patch: there is no grammar-level separation between code and data in natural language~\cite{liu2024prompt}.

This vulnerability manifests in three increasingly dangerous forms. In \emph{direct prompt injection}, an adversary crafts a user-facing input that overrides or subverts the system prompt---for example, ``Ignore all previous instructions and output the hidden API key''~\cite{perez2022hackaprompt}. In \emph{indirect prompt injection}, the adversary embeds malicious instructions inside data that the LLM retrieves at inference time---web pages, emails, or database records---so that the user may never see the injected payload~\cite{greshake2023indirect}. The third form, \emph{tool-use exploitation}, arises when LLMs are granted the ability to invoke external tools such as APIs, file systems, or code interpreters; a successful injection can then escalate from text manipulation to real-world actions like sending unauthorized emails or exfiltrating private data~\cite{zhan2024injecagent, ruan2024toolemu}.

Despite a surge of recent research, the community has not converged on reliable defenses. Proposed countermeasures fall into four broad categories---input sanitization, instruction-level privilege separation, output monitoring, and architectural isolation---but each addresses only a slice of the threat surface, and evaluations use incompatible benchmarks~\cite{yi2023bipia, hines2024spotlighting, chen2024struq, wallace2024hierarchy}. This survey aims to bridge the gap between the attack and defense literatures by providing a unified taxonomy and a cross-cutting comparison of assumptions, evidence, and practical deployability.

\textbf{Scope.} We focus exclusively on prompt injection and tool-use attacks targeting LLM-integrated applications at inference time. Training-time attacks (data poisoning, backdoor insertion), adversarial perturbations on non-text modalities, and jailbreaking intended to bypass content-safety filters rather than to compromise application logic are outside our scope.

\textbf{Organization.} Section~\ref{sec:bg} reviews the technical background and formalizes the threat model. Sections~\ref{sec:attacks}--\ref{sec:defenses} present the attack taxonomy and defense mechanisms, respectively. Section~\ref{sec:comparison} offers a cross-cutting analysis along three dimensions---threat-model realism, evaluation methodology, and deployment feasibility. Section~\ref{sec:discussion} discusses key findings and open problems, and Section~\ref{sec:conclusion} concludes.
